\section{Introduction}

\subsection{Scientific Merit, Importance, and Challenges}

\subsubsection{Scientific Merit}
Many applied data science systems are sequential decision systems (SDSs) that
route limited resources under uncertainty. SDSs are scientifically challenging
because they observe only partial feedback: the chosen action reveals an
outcome, but the unchosen alternatives remain counterfactual.

\subsubsection{Importance}
A policy optimized only for aggregate reward can hide subgroup errors and
compound disparities over time when one group has delayed testing, incomplete
history, lower-quality measurements, or less representative calibration data.
Since the quality of the observed context is often uneven across groups, SDSs
are critical to fairness-aware approaches that mediate compounded disparities.

\subsubsection{Challenges}
This work frames the central challenge as fairness-aware contextual bandit
learning. At each round, a policy observes context, chooses an action, receives
reward only for the chosen action, and updates its future decisions. This
framing guides the scientific questions toward whether informative context and
fairness-aware monitoring reduce disparities when context is missing or
degraded. This focus matters because of the challenges identified in both
target domains. In clinical workflows, false-negative or delayed-escalation
gaps can produce direct harm. In quantum-network routing, future communication
systems will allocate scarce probabilistic resources among competing flows.
Together, these domains expose a central challenge: fairness in success
probability and latency is a service-quality issue that must be addressed.

\subsection{Overview of the Work and Relationship to Prior Work}

\subsubsection{Two Complementary Environment Families}
This work studies two complementary environment families. The first is a
quantum-routing evaluation stack based on hybrid contextual bandits for
entanglement routing and qubit allocation~\cite{garcia2026threat_aware_quantum_routing}.
That stack already exposes uncertainty, adversarial or non-stationary
conditions, legacy testbed conventions, and stateful evaluator artifacts. The
second is an embedded clinical simulation motivated by equitable bioinformatics
and fairness auditing work~\cite{garcia2025equitable_bioinformatics}.

\subsubsection{Clinical and Quantum Routing Analogy}
The clinical environment models the routing of scarce tests, retesting
capacity, escalation decisions, diagnostic attention, and model-assisted triage
across sites and cohorts. Under volatile conditions such as a pandemic, demand,
evidence quality, and resource availability can change rapidly. This makes the
clinical routing setting structurally comparable to quantum routing under noisy
links, uncertain entanglement generation, and non-stationary service
conditions. The two-domain design tests whether fairness mechanisms and
policy-update strategies transfer across fundamentally different applications
with the same partial-feedback allocation structure.

\subsubsection{Contribution Relative to Prior Work}
Across both environments, the central concern is that unequal context quality
can turn partial feedback into repeated allocation harm. Prior bandit methods
primarily optimize aggregate reward, while many fairness analyses evaluate
disparity after decisions have already been made. This study evaluates
performance and fairness as coupled outcomes rather than as separate post-hoc
measurements.

\subsection{Preliminary Results and Current Evidence}

\subsubsection{Selected MAB Spectrum}
Here, our MAB spectrum refers to the non-contextual, contextual, and
informative contextual bandit families evaluated in this work.

\subsubsection{Questions and Findings}
This study asks how context affects performance across our MAB spectrum and
how context quality relates to fairness. The current implementation establishes
the evaluation path needed to answer those questions by producing comparable
clinical and quantum outcome records, model-stratified fairness summaries, and
fairness artifacts that preserve the context, action, allocation, reward, and
group fields needed for analysis. It also asks how fairness-aware mediation
affects overall performance. Current clinical smoke results show that the
embedded clinical path can compare multiple model families through allocation
decisions and EQUITAS audit/mitigation context, with model-level efficiency
summaries already available for the clinical fixture. The remaining questions
ask which MAB family is most likely to amplify disparities and which group is
more likely to be harmed in a non-stationary environment. The present evidence
supports the reporting and validation layer needed for those comparisons, but
full claims about active fairness-adjusted quantum learning require an explicit
policy-update mechanism and full experimental validation.
